\section*{Capítulo \ref{ch:b3:facies-models} --- Modelos de fácies}

\begin{solution}{exo:b3:facies-models:1}
Do mar para a terra: offshore, transição offshore com areias de tempestade,
\omterm{def:b3:facies-models:shoreface}{antepraia}, praia, uma planície costeira com pântanos. Cada ambiente passou a repousar
sobre o vizinho situado mais para o mar: a linha de costa deslocou-se para o mar, numa regressão. Uma
transgressão os teria empilhado no sentido inverso, com a planície costeira na
base e lama de offshore no topo, frequentemente com uma superfície de erosão onde a
\omterm{def:b3:facies-models:shoreface}{antepraia} se deslocava.
\end{solution}

\begin{solution}{exo:b3:facies-models:2}
(a) Preenchimento de canal ou \omterm{def:b3:facies-models:rivers}{barra de pontal}: base erosiva com cascalho, estratificações cruzadas,
depois marcas onduladas e lama com raízes. (b) \omterm{def:b3:facies-models:shoreface}{Antepraia} progradante ou frente deltaica:
lama com camadas de areia de tempestade e \omterm{def:b3:facies-models:hcs}{estratificação cruzada hummocky} na base,
areia de \omterm{def:b3:facies-models:shoreface}{antepraia} com estratificação cruzada e praia acima. (c) Areia de \omterm{def:b3:facies-models:rivers}{rio entrelaçado} empilhada,
duna eólica ou banco carbonático: conjuntos de estratificação cruzada de granulometria
constante, com grandes foresets de duna no caso eólico ou grãos esqueléticos e oolíticos
no caso do banco.
\end{solution}

\begin{solution}{exo:b3:facies-models:3}
$\phi = (2.65 - 2.40)/(2.65 - 1.0) = 0.152$; com salmoura,
$0.25/1.55 = 0.161$.
\end{solution}

\begin{solution}{exo:b3:facies-models:4}
$e^{-2\pi z/\lambda} = e^{-\pi} = 0.043$ em $z = \lambda/2$: cerca de
\qty{38}{m} e \qty{55}{m}. Em \qty{20}{m}: $e^{-2\pi\times20/75} = 0.19$
e $e^{-2\pi\times20/110} = 0.32$ do movimento da superfície.
\end{solution}

\begin{solution}{exo:b3:facies-models:5}
$(60 - 47.6)/(189 - 47.6) = 0.088$; $(60 - 52.6)/(189 - 52.6) = 0.054$. Um
valor de matriz incerto em \qty{5}{\micro\second/ft} altera a porosidade em
um terço: a equação de média temporal exige um valor de matriz verificado em testemunhos.
\end{solution}

\begin{solution}{exo:b3:facies-models:6}
$(6.0 + 3.5 + 8.0 + 2.5)/40 = 20/40 = 0.50$. Volume de poros: $20\times0.22
\times10^6 = \qty{4.4e6}{m^3}$ por quilômetro quadrado.
\end{solution}

\begin{solution}{exo:b3:facies-models:7}
Perto do cânion: somente Ta, com as divisões superiores removidas pelo fluxo seguinte; ao longo do
leque: Tc (marcas onduladas) e Td (silte) sob Te (lama). A corrente perde velocidade
e sua carga grossa com a distância; apenas a fração fina percorre
\qty{30}{km}, e é depositada por um fluxo mais lento que forma marcas onduladas, e não
por decantação de uma suspensão densa.
\end{solution}

\begin{solution}{exo:b3:facies-models:8}
\omterm{def:b3:facies-models:rivers}{Barra de pontal} de um \omterm{def:b3:facies-models:rivers}{rio meandrante}: pavimento residual (\textit{lag}) sobre o fundo erodido do canal,
areia com estratificação cruzada do canal profundo, areia com marcas onduladas no topo da barra,
lama de planície de inundação com raízes. A sucessão registra o canal desde o fundo até
a margem, de modo que o canal tinha cerca de \qty{9}{m} de profundidade (um pouco menos, após
compactação da lama). A base não entra no cálculo: trata-se de uma superfície erosiva, e
o que jazia abaixo pode ter sido qualquer depósito mais antigo.
\end{solution}

\begin{solution}{exo:b3:facies-models:9}
$60/40\,000 = 0.0015$, \ang{0.09}: uma rampa. $300/1500 = 0.20$, \ang{11.3},
após um topo plano de \qty{5}{m} de profundidade: uma \omterm{def:b3:facies-models:ramp}{plataforma com borda}, com grãos de recife ou de banco
na margem, e blocos e fluxos granulares na encosta abaixo.
\end{solution}

\begin{solution}{exo:b3:facies-models:10}
A faixa leva um tempo $w/v$ para passar por um ponto fixo, durante o qual o substrato
sobe em relação ao nível do mar, à taxa $A$: a areia preenche $h = A\,w/v$.
$2000\times0.2/20 = \qty{20}{m}$. Com $A = 0$, a \omterm{def:b3:facies-models:shoreface}{antepraia} deixa apenas
a espessura do próprio perfil, um lençol delgado, pois as faixas deslizam uma sobre
outra sem agradação. A razão $A/v$ é a tangente do ângulo da
trajetória da proposição do Capítulo~\ref{ch:b3:sequence-stratigraphy}:
a espessura de cada fácies é sua largura multiplicada por essa tangente.
\end{solution}

\begin{solution}{exo:b3:facies-models:11}
Densidade: $(2.65 - 2.30)/1.65 = 0.212$; sônico: $(86 - 52.6)/136.4 = 0.245$.
Com $\phi = 0.245$, $\rho_{\mathrm{ma}} = (2.30 - 0.245)/(1 - 0.245) =
\qty{2.72}{g/cm^3}$: mais denso que o quartzo, como num arenito cimentado por
calcita (2.71) ou contendo grãos de calcita.
\end{solution}

\begin{solution}{exo:b3:facies-models:12}
Lama de offshore sobre uma camada de carvão põe em contato ambientes não adjacentes: uma
\omterm{def:b3:sequence-stratigraphy:parasequence}{superfície de inundação}, aqui uma \omterm{def:b3:sequence-stratigraphy:surfaces}{superfície transgressiva} de erosão por ondas, com
detritos de conchas como \textit{lag}. À medida que a linha de costa migrava para o continente, as ondas rebaixaram o
substrato até cerca da \omterm{def:b3:facies-models:shoreface}{base das ondas de tempo bom}: a praia, a \omterm{def:b3:facies-models:shoreface}{antepraia} superior e o
topo da planície costeira, até cerca de \qty{10}{m} de seção, podem estar
ausentes.
\end{solution}

\begin{solution}{pb:b3:facies-models:1}
\textbf{1.} Giz (frio, baixa densidade), folhelho a partir de \qty{766}{m} (quente), calcário a partir de
\qty{866}{m} (frio, denso, rápido), arenito a partir de \qty{1046}{m} (em forma de
sino), xisto a partir de \qty{1166}{m}.

\textbf{2.} Giz \qty{766}{m}, folhelho 100, calcário 180, arenito
\qty{120}{m}; base da cobertura em $441.0 - 1166 = \qty{-725}{m}$.

\textbf{3.} $(2.71 - 2.03)/1.71 = 0.40$ perto da superfície, $(2.71 -
2.36)/1.71 = 0.20$ a \qty{750}{m}: o giz perde metade de sua porosidade
sob o peso da sobrecarga, por compactação.

\textbf{4.} Seus grãos são argilominerais de densidade incerta, e sua água está
parcialmente ligada a eles: as premissas da equação deixam de valer.

\textbf{5.} Embasamento cristalino: a cobertura é atravessada, e nenhum reservatório é
esperado abaixo.

\textbf{6.} $(2.71 - 2.573)/1.71 = 0.080$.

\textbf{7.} $(58.9 - 47.6)/141.4 = 0.080$: o mesmo valor.

\textbf{8.} $(2.65 - 2.32)/1.65 = 0.200$.

\textbf{9.} $(79.9 - 52.6)/136.4 = 0.200$: o mesmo valor.

\textbf{10.} $0.33/1.55 = 0.213$, cerca de um ponto e um terço a mais.

\textbf{11.} O arenito, com \qty{20}{\percent} de porosidade em suas areias;
o folhelho acima do calcário o sela, e as lamas de planície de inundação
dividem o arenito em compartimentos.

\textbf{12.} Uma campânula: raio gama baixo na areia com estratificação cruzada, subindo
pelo intervalo de granodecrescência até os valores altos da lama com raízes, sobre uma base abrupta.

\textbf{13.} Uma \omterm{def:b3:facies-models:association}{sucessão com granodecrescência ascendente}: resíduo de fundo de canal sobre uma
base erosiva, dunas de canal (estratificação cruzada), topo da \omterm{def:b3:facies-models:rivers}{barra de pontal} (marcas onduladas), planície de inundação
(lama com raízes).

\textbf{14.} O canal, a \omterm{def:b3:facies-models:rivers}{barra de pontal} dele na margem interna, diques marginais e uma
planície de inundação com solos, lado a lado.

\textbf{15.} Aproximadamente a espessura de areia e do intervalo de granodecrescência, $10.4 + 3.2 =
\qty{13.6}{m}$, um pouco mais antes da compactação.

\textbf{16.} Um \omterm{def:b3:facies-models:rivers}{rio meandrante} cujo cinturão de canais se deslocou repetidamente por \omterm{def:b3:facies-models:rivers}{avulsão}
enquanto a bacia subsidia, cada novo canal escavando a planície de inundação
deixada pelo anterior.

\textbf{17.} Não: cada base é erosiva; a lama da planície de inundação abaixo dela e
a areia do canal acima dela não precisa ter sido vizinha da outra.

\textbf{18.} $(60 - 30)/(105 - 30) = 0.40$.

\textbf{19.} Mínimo: $14\times(0.60 + 0.40\times0.25) = \qty{9.80}{m}$;
máximo: $16\times(0.65 + 0.40\times0.20) = \qty{11.68}{m}$.

\textbf{20.} Areia limpa \qty{58.8}{m}; intervalos de granodecrescência \qty{33.4}{m}; lama
\qty{27.8}{m}.

\textbf{21.} Líquido $= 58.8 + 0.40\times33.4 = \qty{72.2}{m}$; \omterm{def:b3:facies-models:ntg}{razão net-to-gross}
$72.2/120 = 0.60$.

\textbf{22.} Frações 0.20 e 0.60: \qty{65.5}{m} e \qty{78.8}{m}, uma
dispersão de \qty{13}{m} apenas pela escolha do corte.

\textbf{23.} $72.2\times0.20\times10^6 = \qty{1.44e7}{m^3}$ por
quilômetro quadrado.

\textbf{24.} A espessura líquida de arenito triássico em Varden-2 é
\textbf{\qty[text-series-to-math=true]{72.2}{m}}.
\end{solution}
